%HEADER

%<---!!!!!!!!!!!!!!! MAKRO-DEFIONITIONEN; BITTE NICHT VERAENDERN !!!!!!!!!!
%<--- ARBEIT EINSEITIG
\def\makroEinseitig{
%KOMA-Script-Klasse: scrreprt
%deutsches Design, Schriftgröße 12, DIN A4
%Literaturverzeichnis und Index in Inhaltsverzerzeichnis einbinden
\documentclass[12pt,a4paper,listof=totoc,oneside]{scrreprt}
%Seitenspiegel einstellen
\usepackage[a4paper]{geometry}
\geometry{a4paper,left=30mm,right=25mm,
bottom=20mm,top=15mm,bindingoffset=2mm,
includehead,includefoot}}
% ARBEIT EINSEITIG --->

\def\makroZweiseitig{
%<--- ARBEIT ZWEISEITIG
%KOMA-Script-Klasse: scrreprt
%deutsches Design, zweiseitig
%Literaturverzeichnis und Index in Inhaltsverzerzeichnis einbinden
\documentclass[12pt,a4paper,listof=totoc,twoside, headsepline]{scrreprt}
\usepackage[a4paper]{geometry}
\geometry{a4paper,left=25mm,right=25mm,
bottom=20mm,top=15mm,bindingoffset=2mm,
includehead,includefoot}}
% ARBEIT ZWEISEITIG --->

%<--- Einstellungen Kopfzeile
\def\makroFH-Kopfzeilenstil{
\pagestyle{scrheadings} 
\setheadsepline{0.4pt}
\pagestyle{scrheadings}
\renewcommand*{\chapterpagestyle}{scrheadings}}
%Einstellungen Kopfzeile --->
%!!!!!!!!!!!!!!! MAKRO-DEFIONITIONEN; BITTE NICHT VERAENDERN !!!!!!!!!!--->


%AUSWAHL: TEXT EINSEITIG (ja/nein)
\makroEinseitig
%\makroZweiseitig

%schalte Umlaute frei
\usepackage[english, ngerman]{babel}
%passende Codierung
\usepackage[utf8]{inputenc}
%Seitenspiegel einzustellen
\usepackage[a4paper]{geometry}

%%%%%%%%%%%%%%%%%%%%%%%%%%%%%%%%%%%%%%%%%%%%%%%%%%%%%%%%%%%%%%%%%%%%%%%%%%%%%%%%
% Deutsche und englische Referenzen
\usepackage{babelbib}

% Utopia Schriftart
\usepackage{utopia}

% Farbige Überschriften
% https://www.overleaf.com/learn/latex/Using_colours_in_LaTeX
\usepackage[dvipsnames]{xcolor}
\usepackage{sectsty}
\chapterfont{\color{CornflowerBlue}}
\sectionfont{\color{Cerulean}}
\subsectionfont{\color{Cerulean}}
\subsubsectionfont{\color{Cerulean}}
\paragraphfont{\color{Cerulean}}
%%%%%%%%%%%%%%%%%%%%%%%%%%%%%%%%%%%%%%%%%%%%%%%%%%%%%%%%%%%%%%%%%%%%%%%%%%%%%%%%

%Mathepaket
\usepackage{amsmath}
%Symbole
\usepackage{amssymb}
%griechische Symbole
\usepackage{upgreek}
%weitere Symbole
\usepackage{pxfonts}
% Phonetischen Alphabete für LaTeX
\usepackage{tipa}
%farbige Schriften
%\usepackage{color}
\usepackage{scrhack}
%Bilder fixieren
\usepackage{float}
%Grafiken einbinden
\usepackage{graphicx}
% Kopf- und Fußzeilen
\usepackage[automark,standardstyle,markusedcase]{scrlayer-scrpage}
% deutsche Überschriften
\usepackage[ngerman]{translator}
% Kopfzeilenabstand festlegen
\setlength{\headheight}{10mm}
%Abb. statt Abbildung
\usepackage{caption3}
\addto\captionsngerman{
\renewcommand{\figurename}{Abb.}
\renewcommand{\tablename}{Tab.}
}
%Glossar-Pakage
% HYPERREF VOR GLOSSARIES EINFÜGEN
\usepackage{hyperref}

\usepackage[
nonumberlist, %keine Seitenzahlen anzeigen
acronym,      %ein Abkürzungsverzeichnis erstellen
toc]          %Einträge im Inhaltsverzeichnis      
{glossaries}
\usepackage{cite}

\hypersetup{ colorlinks,
linkcolor=MidnightBlue,
filecolor=MidnightBlue,
urlcolor=MidnightBlue,
citecolor=MidnightBlue }

%Glossar einschalten
\makeglossaries


%fertigen Kopfzeilenstil aktivieren
\makroFH-Kopfzeilenstil

%Zeilenabstand * 1.25 (default)
\renewcommand{\baselinestretch}{1.25}\normalsize
%(Kommentar entfernen, um Zeilenabstand
% auf 1,5-fache Groesse zu ueberschreiben)
%\renewcommand{\baselinestretch}{1.50}\normalsize

% KOMPILATION
% pdflatex Arbeit_Main.tex
% makeglossaries Arbeit_Main
% bibtex Arbeit_Main
% 2 * pdflatex Arbeit_Main.tex
%

\newcommand{\todo}[1]{\textcolor{red}{\textbf{TODO #1}}}

\begin{document}

%Befehle für Symbole
%\newglossaryentry{symb:Pi}{
%name=$\pi$,
%description={Die Kreiszahl.},
%}
%\newglossaryentry{symb:Phi}{
%name=$\varphi$,
%description={Ein beliebiger Winkel.}
%}
%\newglossaryentry{symb:Lambda}{
%name=$\lambda$,
%description={Eine beliebige Zahl, mit der der nachfolgende Ausdruck
%multipliziert wird.}
%}

%Akronyme

%\newacronym{MS}{MS}{Microsoft}

%Eine Akronym mit Glossareintrag
\newacronym{PCA}{PCA}{Primary Component Analysis\protect\glsadd{glos:PCA}}

%Glossareinträge

%\newglossaryentry{glos:AD}{
%name=Active Directory,
%description={Active Directory ist in einem Windows 2000/" "Windows
%Server 2003-Netzwerk der Verzeichnisdienst, der die zentrale ...
%}
%}

\newglossaryentry
{glos:Predictive_Analytics}{name=Predictive Analytics, description=
{Die Analyse von Daten mit Hilfe mathematischer und statistischer Methoden.
Das Ziel von Predictive Analytics ist die Erstellung von Prognosen zur
Unterstützung der Entscheidungsfindung}
}

\newglossaryentry
{glos:PCA}{name=Primary Component Analysis, description=
{Eine Methode zur Variablenreduktion. Dabei wird ein neuer Satz von
Basisvektoren berechnet, der die Varianz der Daten am Besten beschreibt}
}


%Titelseite
% Seitennummer aus
\thispagestyle{empty}
\begin{titlepage}
	\vspace{3cm}

%%\begin{center}
%%	\Huge	
%%	HOCHSCHULE LANDSHUT \\
%%	\Large
%%	FAKULTÄT INFORMATIK
%%\end{center}

%%\vspace{1cm}

\begin{center}
	\includegraphics[scale=0.8]{Grafiken/hl-logo.pdf}  
\end{center}

\vspace{2.5cm}

\begin{center}
  \Large FAKULTÄT INFORMATIK
\end{center}

\vspace{1cm}
\begin{center}
	\Huge
	\textbf{Masterarbeit}\\
\end{center}

\vspace{1cm}

\begin{center}
	\Large
	\textsc{Predictive Analytics im öffentlichen Sektor - Anwendungsmöglichkeiten,
    Chancen und Risiken}\\
\end{center}

\vspace{1.5cm}

\begin{center}
	\Large
	Dmitrij Novikov
\end{center}

\vspace{2cm}
\begin{center}


	\large
	%\the\month\,/\,\the\year
Eingereicht am \today
\end{center}

\vspace{2cm}
\begin{center}
	\large
	Betreuer: Prof. Dr. Wunderlich
\end{center}

\end{titlepage}


%Füge leere Seite ein (optional)
\clearpage 
\ifodd\count0\else 
\thispagestyle{empty} 
\hbox{}\newpage 
\fi 
%Verwende Muster-Erklärung zur selbstständigen Arbeit 
%Erklärung
\thispagestyle{empty}
\vspace{15mm}
\begin{center}
\textbf{\underline{ERKLÄRUNG ZUR MASTERARBEIT}}
\end{center}
\vspace{25mm}
\begin{center}
\large
%Name, Vorname des Studierenden:  \\
\large
Novikov, Dmitrij
\end{center}
\vspace{25mm}

\begin{center}
\huge
Hochschule Landshut \\
Fakultät Informatik 
\end{center}
\vspace{10mm}

\begin{center}
\large
Hiermit erkläre ich, dass ich die Arbeit selbständig 
verfasst, noch nicht anderweitig für Prüfungszwecke 
vorgelegt, keine anderen als die angegebenen Quellen 
oder Hilfsmittel benützt, sowie wörtliche und sinngemäße 
Zitate als solche gekennzeichnet habe.  \\
\end{center}
\vspace{55mm}

\begin{center}
....................\hspace{40mm}....................................................\\

(Datum)\hspace{47mm}(Unterschrift des Studierenden)
\end{center}

%Füge leere Seite ein (optional)
\clearpage 
\ifodd\count0\else 
\thispagestyle{empty} 
\hbox{}\newpage 
\fi 

\begin{abstract}
\begin{center}
\Huge
\emph{\textbf{Abstract}}
\end{center}
\normalsize
\vspace{15mm}
\textit{\todo{Abstract}}
\end{abstract}


\tableofcontents
\setcounter{page}{1}

\chapter{Grundlagen von Predictive Analytics}

Textkram Textkram Textkram Textkram Textkram Textkram Textkram \\ \\
Mehr Textkram Mehr Textkram Mehr Textkram Mehr Textkram Mehr Textkram Mehr
Textkram \\ \\
Abkürzung (Akronym) \gls{MS} und Glossareintrag \gls{glos:AntwD}. \\ \\
Zitieren auf Deutsch \cite{ART12d} und Englisch \cite{ART12e}. \\ \\
Hier ein Link \url{www.kram.de}.

\cite{Thapa_Parycek} und \cite{Mergel}

\section{Motivation - Die Schwierigkeit von Vorhersagen}

\section{Begriffsdefinition und -Abgrenzung}

\section{Bestandteile von Predictive Analytics}

\subsection{CRISP-DM Vorgehensmodell}

\subsection{Daten}

\subsubsection{Information Management (Data Warehouse)}

\subsubsection{Datenabhängige Ziele}

\paragraph{\ldots}

\paragraph{Beschreibung der Daten}

\paragraph{Klassifikation}

\paragraph{Zeitreihenanalyse}

\subsection{Methoden}

\subsubsection{Die Großen Drei}

\paragraph{Regression}

\paragraph{Entscheidungsbäume}

\paragraph{Neuronale Netze}

\subsubsection{Zeitreihenanalyse}

\subsubsection{\ldots}

\subsection{Werkzeuge}

\subsection{(Kreative) Freiheitsgrade bei der Implementierung eines
  Vorhersagemodells}

\section{Anwendungsbeispiele}

\subsection{Allgemeine Anwendungsfelder (siehe SAP Artikel)}

\subsection{\ldots}

\section{Allgemeine Risiken bei der Nutzung von Predictive Analytics}

\subsection{Risiko der Unverhältnismäßigkeit}

\subsection{Risiken bei der Datenerhebung}

\subsubsection{\ldots}

\subsection{Risiken bei der Interpretation der Daten}

\subsection{Risiken bei der Interpretation der Ergebnisse}

\subsubsection{Prognose verändert das Verhalten des Systems 
%  (``Selffulfilling Prophecy'')}
  (\glqq{Selffulfilling} Prophecy\grqq})

\subsubsection{Ignorieren der Unsicherheiten bei der Prognose}

\chapter{Predictive Analytics im öffentlichen Sektor}

\section{Übersicht}

\section{Thematische Schwerpunkte}
\todo{Wo könnte der größte Nutzen liegen?}

\todo{vorläufig:}

\subsection{Öffentliche Verwaltung}

\subsubsection{Daten in der öffentlichen Verwaltung}

\paragraph{Open Data Portale ( GovData.de )}

\subsubsection{Organisationseinheiten für Predictive Analytics}

\paragraph{\ldots}

\subsection{Gesundheit}

\subsection{Bildung}

\chapter{Gegenwärtige Anwendungen}

\section{Predictive Analytics in Finanzbehörden}

\section{Google Flu Trends}

\section{Predictive Policing}

\section{Cambridge Analytica}

% TODO Auflistungstabelle
% Kategorie | Anwendung | Beispiel

\chapter{Zukunftsperspektiven}

\section{\ldots}

\section{Einschätzung der politischen Stimmungslage}

\section{Vorhersage kultureller Unterschiede}

\chapter{Zusammenfassung und Fazit}

%%%%%%%%%%%%%%%%%%%%%%%%%%%%%%%%%%%%%%%%%%%%%%%%%%%%%%%%%%%%%%%%%%%%%%%%%%%%%%%%
\appendix
\chapter{Probablility Scoring - Vertiefung}

\chapter{Predictive Analytics mit R}

\section{Tabellenkalkulationen versus R}

Inhalt vom ersten Anhang.

\section{Beispielhafte Auswertung eines Datensatzes}

Inhalt des zweiten Anhang.

%%%%%%%%%%%%%%%%%%%%%%%%%%%%%%%%%%%%%%%%%%%%%%%%%%%%%%%%%%%%%%%%%%%%%%%%%%%%%%%%
\bibliography{Bibliographie/Bibli.bib}{}
\bibliographystyle{babplain}

\listoffigures
\listoftables
\printglossary[title=Glossar]
\printglossary[type=\acronymtype, title=Akronyme]
% keine deutschen Überschriften
%\printglossaries 

\end{document}


